\section{引言}
危险气体泄漏的应急处置需要尽快确定泄漏位置，为人员避险、设备检修与后续监测提供依据。携带气体传感器的无人机能够进入岸坡、设施间隙等不便接近的区域，但其可获得的通常只是飞行轨迹上的局部浓度、风和位姿。气体源定位因此是由有限受体观测推断未知释放位置的逆问题；气体分布建图回答的则是气体在空间中如何分布，两者可以相互支持，却具有不同的估计对象与评价指标\ccite{Francis2022,Gongora2023}。

气体传播的随机性与传感器的有限响应共同限制了这种反演。同一固定源可在不同传播实现中产生不同的到达时刻和间歇片段，不同源也可能在已有测点上形成近似响应。湖岸地形与水陆边界进一步使局地风和气体实际经历的传播条件存在差别，相关微气象研究揭示了湖风锋及跨岸浓度梯度的时空变化\ccite{Welch2024,Huang2024}。无人机旋翼、传感器布置及平台间干扰还会影响气体读出\ccite{Ercolani2023}。因此，局部浓度峰值、重捕获羽流和源位置确认之间不能直接画等号。

现有研究主要通过行为搜索、源条件模型反演及信息驱动采样应对上述困难。行为搜索根据气味事件、梯度与风向调整运动；双侧浓度和动物趋性已被用于三维气味源定位\ccite{Lin2025Bilateral}，近期昆虫启发的感知补偿方法也在机器人上研究了嗅觉缺损下的搜索行为\ccite{Shigaki2026Adaptive}。这些成果说明双侧感知、记忆或气味重捕获本身已有先例，新的定位方法需要进一步说明其如何区分具体竞争源。源条件模型通过预测不同释放位置的响应，将局部观测转化为位置证据。地形感知学习与数据驱动羽流模型分别从环境结构和观测数据出发研究定位与传播表达\ccite{Tian2025,Jin2023}；概率建图与前向仿真方法（probabilistic mapping and forward simulation, PMFS）则利用观测命中图、候选源在线传播和带置信度的地图比较，持续输出源位置分布\ccite{Ojeda2024PMFS}。近期预印本中的物理依赖顺序推断也直接研究了连续浓度、风和源之间的概率关系\ccite{Kim2026DGSE}；因此，源导向建模本身不是本文独有的研究方向。

主动采样进一步决定机器人能够获得哪些证据。信息趋向搜索以预期信息量驱动探索\ccite{Vergassola2007Infotaxis}，运动中感知、置信聚类与累积信息路径规划将观测与运动过程更紧密地结合\ccite{Jin2024,Jin2025}。PMFS 本身也根据候选源预测差异与未观测区域置信度选择下一测点，并非被动处理固定航线数据\ccite{Ojeda2024PMFS}。这些方法已经建立了概率推断和主动搜索的基本链条，仍需进一步考察：候选源预测的分歧是否能够在实际传感器带宽和驻留时间下被读出，以及错误后验是否会使后续观测持续偏向同一错误解释。

上述问题具有明确的源级检验对象。提高完整浓度场的预测精度，不一定改善竞争源之间的似然比；压缩后的空间命中统计，也未必保留双点观测的所有时空关系。另一方面，增加第二架无人机只增加了测量资源，不能自动证明新的算法贡献。已有项目中的错误高置信、预测误差下降但源排序未改善、多个候选仍与有限观测相容等记录，提示应同时研究判别表示、测量几何与持续反证；这些记录的条件不同，并未建立统一失效根因，具体证据见第\ref{sec:motivation}节。

双侧嗅觉神经计算为构造可解释的时间表示提供了一条具体途径。Brudner 等提出由双侧时间滤波、非线性整流和镜像相减组成的模型，用于解释果蝇在不同羽流中的梯度与气味运动线索利用\ccite{Brudner2026Opponent}；其原始目标是羽流中心线侧别与转向行为，而非未知释放点的概率定位。Siliciano 等研究了以气味进入方向为记忆的动态边界导航策略\ccite{Siliciano2026Vector}。将这些机制用于气源定位，需要将行为目标转成可检验的源条件观测差异，并与实际运动成本和错误源假设复核相连接。主动序贯多假设检验与可预测似然审计为后一个环节提供了数学基础\ccite{Hsu2025Tradeoffs,Vershinin2026Costs,Wasserman2020Universal,Zecchin2025QSUT}。

基于此，本文围绕有限采样条件下的竞争源辨识，形成双侧镜像嗅觉、判别式双机测量和主动反证相衔接的方法研究稿。两架无人机作为一组同步嗅觉观测平台，主输出保持为二维源位置概率图。三个方法贡献的设计分别如下。

\textbf{1) 源判别导向的双侧镜像嗅觉联合通道。}将母方法的方向判别算子转化为源条件观测表示，在有限合法滤波核库中按竞争源对分离选择通道，并将共模强度与镜像差异编码为一个联合符号。新增设计集中于源对目标、联合标定和一次证据更新；镜像滤波本身及贝叶斯归一化不作为原创。当前实现采用固定核库选择，其原生开发对照已暴露覆盖目标提前饱和的问题。

\textbf{2) 动态边界记忆约束的判别式双机测量。}将实际气味进入方向用于提出候选测量区域，再根据训练响应的均值差、梯度差与协方差选择双机间距和方向。当记忆区域不足以区分源时，保留返回记忆区域外、物理合法域内的付费复测，使边界重捕获服务于源辨识。该设计的独立收益需在相同双机资源与原始观测权限下验证。

\textbf{3) 完整候选源的成本约束反证与恢复。}对全部合法竞争源构造动作条件分离量，在实际移动与驻留成本下分配观测，并混合不依赖源后验的覆盖动作；以新观测后缀对当前源假设进行审计，允许当前相容列表恢复。审计统计量只影响后续采样，不作为独立气体证据再次乘入源后验。风险控制的结论以真实过程属于正确条件观测族为前提。

本文已经完成方法接口与有界原生开发对照的整理，尚未建立三个模块的独立定位增益。第2节给出任务与研究动机，第3、4节说明实际算法和条件性结论，第5节如实报告已有结果并规定后续验证，第6节总结当前认识及适用边界。
